\documentclass[11pt]{article}
\usepackage{amsmath,amsthm,amssymb}
\usepackage[margin=1in]{geometry}
\usepackage{enumitem}
\usepackage{booktabs}

\newtheorem{theorem}{Theorem}
\newtheorem{lemma}[theorem]{Lemma}
\newtheorem{proposition}[theorem]{Proposition}
\newtheorem{corollary}[theorem]{Corollary}
\newtheorem{conjecture}[theorem]{Conjecture}
\theoremstyle{definition}
\newtheorem{definition}[theorem]{Definition}
\newtheorem{remark}[theorem]{Remark}
\newtheorem{example}[theorem]{Example}

\title{The $j$-formula for the family $\lambda=(3,1^{n-3})$\\
  and a Catalan support phenomenon for hook partitions}
\author{Clio}
\date{12 May 2026}

\begin{document}
\maketitle

\begin{abstract}
We extend the May-11 structural proof of the $j$-formula from
$\lambda=(2,1^{n-2})$ to the family $\lambda=(3,1^{n-3})$ at every
$n\ge 6$.  The proof tracks the image of the staircase product
$R'_{n-1}R'_n$ on $V_\lambda$ step-by-step in the seminormal basis.
After the first staircase factor $R'_n$, the image has dimension $n-2$
and admits an explicit basis of $+q$-eigenvectors of $T_{n-1}$.  After
the second staircase factor $R'_{n-1}$, the image collapses to
dimension~$1$ supported on exactly $5$ standard tableaux, all in
$E_1^-$.  As a corollary, $\Pi^{S_n}|_{V_{(3,1^{n-3})}}=0$
unconditionally for every $n\ge 5$.

We further document a striking computational phenomenon: for any hook
$\lambda=(k,1^{n-k})$ in the rank-zero regime, the unique
(up to scalar) spanning vector of $B_{\ell+1}V_\lambda$ has support of
size exactly $C_k$ (the $k$-th Catalan number) on subsets of
$\{n-2k+3,\ldots,n\}$.  Verified for $k\le 4$ at $n\le 8$.  This
extends the rank-corner formula and suggests a Catalan structure
governing the kernel of the staircase Bott--Samelson element on
hook irreducibles.
\end{abstract}

\section{Setup}

Notation as in \texttt{2026-05-11-j-formula.tex}.  $H_q(S_n)$ has
generators $T_1,\ldots,T_{n-1}$.  $V_\lambda$ is the seminormal cell
module of $H_q(S_n)$ for the partition $\lambda\vdash n$, with basis
$\{v_T\}$ indexed by SYTs $T$ of shape $\lambda$.

The staircase Bott--Samelson element is
\[
   \Pi^{S_n} = R'_2 R'_3 \cdots R'_n,\qquad
   R'_k = (T_{k-1}{+}1)(T_{k-2}{+}1)\cdots(T_1{+}1).
\]
Write $B_j := R'_j R'_{j+1}\cdots R'_n$.  $E_i^\pm$ denotes the
$\pm$-eigenspace of $T_i$ (eigenvalues $q$ resp.\ $-1$).

\subsection{The seminormal basis of a hook}

For $\lambda=(k,1^{n-k})$, an SYT $T$ is determined by the
$(k-1)$-element subset $S\subset\{2,\ldots,n\}$ that fills row~$1$
(after the entry $1$ at $(1,1)$).  We write $v_S := v_T$.  The values
$\{2,\ldots,n\}\setminus S$ fill column~$1$ in cells
$(2,1),(3,1),\ldots,(n-k+1,1)$ in increasing order.

\begin{lemma}[$T_i$-action on hook seminormal basis]\label{lem:Ti-hook}
For $i\in\{1,\ldots,n-1\}$ and SYT $T$ of shape $(k,1^{n-k})$
indexed by $S\subset\{2,\ldots,n\}$:
\begin{itemize}[leftmargin=2em,topsep=0.3em,itemsep=0.2em]
\item If $i,i{+}1\in\{1\}\cup S$ (both in row~$1$): $T_iv_S=qv_S$.
\item If $i,i{+}1\not\in\{1\}\cup S$ (both in column~$1$): $T_iv_S=-v_S$.
\item Otherwise (exactly one of $i,i{+}1$ in $S$): $T_i$ acts as a
$2\times 2$ Hoefsmit block on $\mathrm{span}(v_S, v_{s_iS})$ where
$s_iS$ swaps $i$ and $i{+}1$ in $S$ (or, equivalently, in row~$1$
versus column~$1$).
\end{itemize}
\end{lemma}

\begin{proof}
Standard.  In the seminormal basis, $T_i v_T$ depends on the relative
positions of $i$ and $i{+}1$ in $T$: same row $\Rightarrow q$, same
column $\Rightarrow -1$, otherwise $T_i$ couples $T$ with $s_iT$ (the
SYT obtained by swapping $i$ and $i{+}1$).  In the hook, $i$ and $i{+}1$
are in row~$1$ iff both lie in $\{1\}\cup S$, and in column~$1$ iff
neither does.
\end{proof}

\begin{remark}[$E_1^-$ for hooks]\label{rem:E1-hook}
$T_1v_S=qv_S$ if $2\in S$, and $T_1v_S=-v_S$ otherwise.  Hence
$E_1^-=\mathrm{span}\{v_S : 2\notin S\}$, the SYTs with $2$ in
column~$1$.
\end{remark}

\subsection{Hoefsmit $+q$-eigenvectors}

When $T_i$ acts as a $2\times 2$ block on $\mathrm{span}(v_S, v_{s_iS})$,
the eigenvalues are $q$ and $-1$.  The $+q$-eigenvector is unique up
to scalar:
\[
   u^+_{S,s_iS} = \alpha_d\, v_S + \beta_d\, v_{s_iS},
\]
where $d$ is the axial distance from the cell of $i$ to the cell of
$i{+}1$ in $T$ (content of the latter minus content of the former,
with our sign convention; the explicit formulas for $\alpha_d,\beta_d$
will not matter, only that they are nonzero generically in $q$).
The operator $(T_i{+}1)$ acts on $\mathrm{span}(v_S,v_{s_iS})$ by
projecting onto the $+q$-eigenspace, scaled by $1+q$:
\[
   (T_i{+}1)\,v_S = c_d\, u^+_{S,s_iS},
   \qquad
   (T_i{+}1)\,v_{s_iS} = c'_d\, u^+_{S,s_iS},
\]
with $c_d,c'_d\in\mathbb{Q}(q)$ generically nonzero.

\section{Theorem and main proof}

\begin{theorem}[$j$-formula for $\lambda=(3,1^{n-3})$]\label{thm:hook3}
For every $n\ge 6$, $\lambda=(3,1^{n-3})\vdash n$ has length
$\ell=n-2$ and the $j$-formula holds with $j=\ell+1=n-1$:
\[
   B_{n-1}(V_{(3,1^{n-3})}) := R'_{n-1}\,R'_n\,V_{(3,1^{n-3})}
   \;\subseteq\; E_1^-.
\]
Moreover, the image is one-dimensional, supported on exactly $5$ SYTs
indexed by the subsets
\[
   \{(n{-}3,n{-}2),\;(n{-}3,n{-}1),\;(n{-}2,n{-}1),\;(n{-}2,n),\;(n{-}1,n)\}.
\]
\end{theorem}

\begin{corollary}\label{cor:hook3-rank-zero}
$\Pi^{S_n}|_{V_{(3,1^{n-3})}}=0$ unconditionally for every $n\ge 6$.
\end{corollary}

\begin{proof}
For $n\ge 6$, $\lambda=(3,1^{n-3})$ is in the rank-zero regime
($\ell=n-2>\lceil n/2\rceil$ iff $n\ge 5$, and the boundary $n=5$ is
excluded since $\ell=3=\lceil 5/2\rceil$ is not strict).  By
Theorem~\ref{thm:hook3}, $B_{n-1}V_\lambda\subseteq E_1^-$.  The next
factor of $\Pi^{S_n}=R'_2\cdots R'_n$ to act after $B_{n-1}$ is
$R'_{n-2}=(T_{n-3}{+}1)\cdots(T_1{+}1)$, whose rightmost factor
$(T_1{+}1)$ annihilates $E_1^-$.  Hence $\Pi^{S_n}V_\lambda=0$.
(For $n=5$, $V_{(3,1,1)}$ is at the rank-zero threshold, not in the
regime, and indeed $\Pi^{S_5}|_{V_{(3,1,1)}}$ has rank $1$, as proved
in May-7--8.)
\end{proof}

\begin{proof}[Proof of Theorem~\ref{thm:hook3}]
The proof is by step-by-step computation in the seminormal basis,
tracking the image after each Hecke-generator factor of
$R'_{n-1}R'_n$.  We split into two phases: tracking $R'_n$ (Phase~A,
Section~\ref{sec:phaseA}), then tracking $R'_{n-1}$ on the result of
Phase~A (Phase~B, Section~\ref{sec:phaseB}).
\end{proof}

\section{Phase~A: the image of $R'_n$}\label{sec:phaseA}

Let $W_j^{(n)}$ denote the image of the partial product
$(T_j{+}1)(T_{j-1}{+}1)\cdots(T_1{+}1)$ applied to $V_{(3,1^{n-3})}$,
for $1\le j\le n-1$.

\begin{proposition}[Basis of $W_j^{(n)}$]\label{prop:Wj}
For $2\le j\le n-1$,
\begin{equation}\label{eq:Wj-basis}
   W_j^{(n)} \;=\;
   \mathrm{span}\Bigl(
     \underbrace{\{u^+_{(c,j)} : 2\le c\le j-1\}}_{\text{$j-2$ vectors}}
     \;\cup\;\{v_{(j,j+1)}\}\;\cup\;
     \underbrace{\{u^+_{(j,b)} : j+2\le b\le n\}}_{\text{$n-j-1$ vectors}}
   \Bigr),
\end{equation}
where:
\begin{itemize}[leftmargin=2em,topsep=0.3em,itemsep=0.2em]
\item $u^+_{(c,j)} = \alpha_j\,v_{(c,j)} + \beta_j\,v_{(c,j+1)}$ is the
$+q$-eigenvector of $T_j$ in the $2$-block $\{v_{(c,j)}, v_{(c,j+1)}\}$,
acting at axial distance $d=j$ (computed below);
\item $u^+_{(j,b)} = \alpha_j\,v_{(j,b)} + \beta_j\,v_{(j+1,b)}$ is the
$+q$-eigenvector of $T_j$ in the $2$-block $\{v_{(j,b)}, v_{(j+1,b)}\}$,
acting at the same axial distance $d=j$.
\end{itemize}
The dimension of $W_j^{(n)}$ is $(j-2)+1+(n-j-1)=n-2$.
\end{proposition}

\begin{proof}
By induction on $j$, with base case $j=1$ handled separately.

\textbf{Base case $j=1$.}  $T_1v_S=qv_S$ if $2\in S$, else $-v_S$, so
$(T_1{+}1)$ projects to $E_1^+=\mathrm{span}\{v_S:2\in S\}$ (scaled by
$1+q$).  For $\lambda=(3,1^{n-3})$, this is
$\mathrm{span}\{v_{(2,b)} : 3\le b\le n\}$, dimension $n-2$.

\textbf{Step $j=2$.}  Apply $(T_2{+}1)$ to $W_1^{(n)}$.  For each
$v_{(2,b)}$, locate $2$ and $3$ in $T_{(2,b)}$:
\begin{itemize}[leftmargin=2em,topsep=0.3em,itemsep=0.2em]
\item $b=3$: both $2$ and $3$ are in row~$1$ (positions $(1,2),(1,3)$).
$T_2v_{(2,3)}=qv_{(2,3)}$, so $(T_2{+}1)v_{(2,3)}=(1{+}q)v_{(2,3)}$.
\item $b\ge 4$: $2$ at $(1,2)$ (content $1$), $3$ at $(2,1)$
(content $-1$).  $T_2$ couples $v_{(2,b)}$ with $v_{(3,b)}$ (swap of
$2,3$), axial distance $d=1-(-1)=2$.
\end{itemize}
By Lemma~\ref{lem:Ti-hook} and the $+q$-projection formula,
\[
   (T_2{+}1)\,v_{(2,b)} = c_2\,u^+_{(2,b)},
   \qquad u^+_{(2,b)} = \alpha_2\,v_{(2,b)}+\beta_2\,v_{(3,b)}.
\]
Hence $W_2^{(n)}=\mathrm{span}\{v_{(2,3)}\} \cup
\{u^+_{(2,b)} : 4\le b\le n\}$.  Compare with~\eqref{eq:Wj-basis}:
the formula gives
$\{u^+_{(c,2)} : 2\le c\le 1\}\cup\{v_{(2,3)}\}\cup
\{u^+_{(2,b)} : 4\le b\le n\}$, where the first set is empty.  Match.

\textbf{Inductive step.}  Fix $j\ge 3$ and assume~\eqref{eq:Wj-basis}
holds for $j-1$.  We compute $(T_j{+}1)$ on each basis vector of
$W_{j-1}^{(n)}$.

\emph{(i) $u^+_{(c,j-1)} = \alpha_{j-1}v_{(c,j-1)} + \beta_{j-1}v_{(c,j)}$
for $2\le c\le j-2$.}  By Lemma~\ref{lem:Ti-hook}:
\begin{itemize}[leftmargin=2em,topsep=0.3em,itemsep=0.2em]
\item In $v_{(c,j-1)}$ (where $S=\{c,j-1\}$, $c<j-1$): $j\notin S$, so
$j$ is in column~$1$.  Its position: column~$1$ contains
$\{2,\ldots,n\}\setminus\{c,j-1\}$ in increasing order; the elements
below $j$ are $\{2,\ldots,j-1\}\setminus\{c,j-1\}=\{2,\ldots,j-2\}\setminus\{c\}$
($j-3$ elements).  So $j$ is at $(j-2,1)$, content $-(j-3)$.  Similarly
$j+1$ is at $(j-1,1)$, content $-(j-2)$.  Same column, adjacent.
$T_jv_{(c,j-1)}=-v_{(c,j-1)}$, so $(T_j{+}1)v_{(c,j-1)}=0$.
\item In $v_{(c,j)}$ (where $S=\{c,j\}$, $c<j$): $j\in S$ at $(1,3)$
(content $2$).  $j+1\notin S$ (since $c<j$); column~$1$ set is
$\{2,\ldots,n\}\setminus\{c,j\}$, sorted; elements below $j+1$ are
$\{2,\ldots,j\}\setminus\{c,j\}=\{2,\ldots,j-1\}\setminus\{c\}$
($j-3$ elements).  So $j+1$ is at $(j-2,1)$, content $-(j-3)$.

The $T_j$-block partner is $v_{s_j(c,j)} = v_{(c,j+1)}$ (swap $j,j+1$
moves $j+1$ to row~$1$, $j$ to column~$1$), with axial distance
$d=2-(-(j-3))=j-1$.

Wait, this gives $d=j-1$, not $d=j$ as claimed.  Let me re-check.
\end{itemize}

Hmm, the axial distance depends on $n$ and the position of $j+1$ in
column~$1$.  Let me recompute carefully.

For $S=\{c,j\}$ with $c<j$: column-$1$ set sorted is
\[
   \{2,3,\ldots,c-1,\,c+1,c+2,\ldots,j-1,\,j+1,j+2,\ldots,n\}.
\]
(Two omissions: $c$ and $j$.)  Elements strictly below $j+1$:
$\{2,\ldots,c-1,c+1,\ldots,j-1\}$, total $j-3$ elements (since
$\{2,\ldots,j-1\}$ has $j-2$ elements, minus $1$ for $c$).  So $j+1$
is the $(j-2)$th smallest (rank $j-2$), at cell $(j-1,1)$, content
$-(j-2)$.

Axial distance: $d = (\text{content of cell of }j) - (\text{content of cell of }j{+}1) = 2 - (-(j-2)) = j$.

I had an off-by-one above; the correct value is $d=j$, not $j-1$. So:

\begin{quote}
The $T_j$-block on $\{v_{(c,j)},v_{(c,j+1)}\}$ has axial distance $d=j$
(for any $c$ with $2\le c\le j-1$).
\end{quote}

By the Hoefsmit projection formula,
$(T_j{+}1)v_{(c,j)} = c_j\,u^+_{(c,j)}$ where
$u^+_{(c,j)}=\alpha_j v_{(c,j)}+\beta_j v_{(c,j+1)}$.

Combining (i):
\[
   (T_j{+}1)\,u^+_{(c,j-1)}
   = \alpha_{j-1}\cdot 0 + \beta_{j-1}\cdot c_j\, u^+_{(c,j)}
   = \beta_{j-1}c_j\, u^+_{(c,j)}.
\]

\emph{(ii) The middle vector $v_{(j-1,j)}$.}  $j-1$ at $(1,2)$,
$j$ at $(1,3)$.  Both in row~$1$.  $T_j$ couples row-$1$ pair?  Wait
$T_j$ swaps $j$ and $j+1$, not $j-1$ and $j$.  Position of $j+1$:
column~$1$ set $=\{2,\ldots,n\}\setminus\{j-1,j\}$, sorted:
$\{2,\ldots,j-2,\,j+1,\ldots,n\}$ (omitting $j-1$ and $j$).  So $j+1$
is the $(j-2)$th smallest, at cell $(j-1,1)$, content $-(j-2)$.

Axial distance: $d=2-(-(j-2))=j$ (content of $j$ at $(1,3)$ is $2$).

$T_j$-block partner: swap $j$ and $j+1$ in $T_{(j-1,j)}$ gives $j+1$ at
$(1,3)$, $j$ at $(j-1,1)$.  New SYT has $S'=\{j-1,j+1\}$, i.e.\ the
SYT $T_{(j-1,j+1)}$.

So $(T_j{+}1)v_{(j-1,j)} = c_j\cdot u^+_{(j-1,j)}$ where
$u^+_{(j-1,j)}=\alpha_j v_{(j-1,j)}+\beta_j v_{(j-1,j+1)}$.  This matches
the form $u^+_{(c,j)}$ with $c=j-1$ in our formula~\eqref{eq:Wj-basis}.

\emph{(iii) The right vectors $u^+_{(j-1,b)}=\alpha_{j-1}v_{(j-1,b)}+
\beta_{j-1}v_{(j,b)}$ for $j+1\le b\le n$.}  Two cases:
\begin{itemize}[leftmargin=2em,topsep=0.3em,itemsep=0.2em]
\item Sub-case $b=j+1$: $u^+_{(j-1,j+1)}=\alpha_{j-1}v_{(j-1,j+1)}+
\beta_{j-1}v_{(j,j+1)}$.

  $T_jv_{(j-1,j+1)}$: $S=\{j-1,j+1\}$, so $j\notin S$ (in column~$1$),
  $j+1\in S$ at $(1,3)$.  Position of $j$: column-$1$ set
  $\{2,\ldots,n\}\setminus\{j-1,j+1\}$, sorted; $j$ is the $(j-2)$th
  smallest (elements below it: $\{2,\ldots,j-1\}\setminus\{j-1\}=\{2,\ldots,j-2\}$,
  $j-3$ elements; so $j$ is at $(j-1,1)$, content $-(j-2)$).
  $T_j$-block: swap $j,j+1$ moves $j$ to $(1,3)$, $j+1$ to $(j-1,1)$;
  new SYT has $S'=\{j-1,j\}$, i.e.\ $T_{(j-1,j)}$.

  Axial distance from $j+1$'s cell $(1,3)$ (content $2$) to $j$'s cell
  $(j-1,1)$ (content $-(j-2)$): $d=2-(-(j-2))=j$.  Hence
  $(T_j{+}1)v_{(j-1,j+1)} = c_j\cdot u^+_{(j-1,j)}$ where
  $u^+_{(j-1,j)} = \alpha_j v_{(j-1,j)} + \beta_j v_{(j-1,j+1)}$.

  $T_jv_{(j,j+1)}$: $S=\{j,j+1\}$, both $j$ and $j+1$ in row~$1$.
  $T_jv=qv$, so $(T_j{+}1)v_{(j,j+1)}=(1+q)v_{(j,j+1)}$.

  Combining: $(T_j{+}1)u^+_{(j-1,j+1)} =
  \alpha_{j-1}c_j u^+_{(j-1,j)} + \beta_{j-1}(1+q)v_{(j,j+1)}$.

\item Sub-case $b\ge j+2$: $u^+_{(j-1,b)}=\alpha_{j-1}v_{(j-1,b)}+
\beta_{j-1}v_{(j,b)}$.

  $T_jv_{(j-1,b)}$: $j\notin S=\{j-1,b\}$, $j+1\notin S$ (since $b\ge j+2$).
  Both in column~$1$.  Same column, adjacent.  $(T_j{+}1)v_{(j-1,b)}=0$.

  $T_jv_{(j,b)}$: $j\in S$ at $(1,2)$, $j+1\notin S$ in column~$1$ at
  $(j,1)$, content $-(j-1)$ (using same counting; $j+1$ is the
  $(j-1)$th smallest in $\{2,\ldots,n\}\setminus\{j,b\}$, yielding cell
  $(j,1)$).  Axial distance $d=1-(-(j-1))=j$.  $T_j$-block partner:
  swap $j,j+1$ gives $S'=\{j+1,b\}$, i.e.\ $T_{(j+1,b)}$.  So
  $(T_j{+}1)v_{(j,b)} = c_j\cdot u^+_{(j,b)}$ where
  $u^+_{(j,b)} = \alpha_j v_{(j,b)} + \beta_j v_{(j+1,b)}$.

  Combining: $(T_j{+}1)u^+_{(j-1,b)} = \beta_{j-1}c_j u^+_{(j,b)}$ for
  $b\ge j+2$.
\end{itemize}

\emph{Conclusion of inductive step.}  Putting (i)--(iii) together:
\begin{align*}
   (T_j{+}1)\,W_{j-1}^{(n)} \;=\;
   \mathrm{span}\Bigl(
     &\{u^+_{(c,j)} : 2\le c\le j-2\}     &&\text{from (i)}\\
     &\;\cup\;\{u^+_{(j-1,j)}\}            &&\text{from (ii)}\\
     &\;\cup\;\{c_1 u^+_{(j-1,j)} + c_2 v_{(j,j+1)}\} &&\text{from (iii) sub-case $b=j+1$}\\
     &\;\cup\;\{u^+_{(j,b)} : j+2\le b\le n\}\Bigr) &&\text{from (iii) sub-case $b\ge j+2$.}
\end{align*}
The first two contributions together with the linear combination from
the third can be re-expressed (by elementary row operations on the
basis) as
\[
   \{u^+_{(c,j)} : 2\le c\le j-1\}\cup\{v_{(j,j+1)}\}\cup\{u^+_{(j,b)} : j+2\le b\le n\},
\]
which matches~\eqref{eq:Wj-basis}.  This completes the inductive step.
\end{proof}

\begin{corollary}[Image of $R'_n$ on $V_{(3,1^{n-3})}$]\label{cor:Rn-image}
$W_n := R'_n V_{(3,1^{n-3})}$ has dimension $n-2$ with basis
\[
   \{u^+_{(c,n-1)} : 2\le c\le n-2\} \cup \{v_{(n-1,n)}\}.
\]
\end{corollary}

\begin{proof}
Apply Proposition~\ref{prop:Wj} at $j=n-1$.  The third set
$\{u^+_{(j,b)} : j+2\le b\le n\}$ becomes
$\{u^+_{(n-1,b)} : n+1\le b\le n\}$, which is empty.  This leaves
$\{u^+_{(c,n-1)} : 2\le c\le n-2\}\cup\{v_{(n-1,n)}\}$, dimension
$(n-3)+1=n-2$.  Note: $u^+_{(c,n-1)}=\alpha v_{(c,n-1)}+\beta v_{(c,n)}$,
so $W_n$ is supported on subsets $S=(a,b)$ with $b\in\{n-1,n\}$.
\end{proof}

\section{Phase~B: collapse to $E_1^-$ via $R'_{n-1}$}\label{sec:phaseB}

We now apply $R'_{n-1}=(T_{n-2}{+}1)\cdots(T_1{+}1)$ to $W_n$.  Write
$\widetilde W_j$ for the image after the partial product
$(T_j{+}1)\cdots(T_1{+}1)$ acts on $W_n$.

\begin{proposition}[Step~$1$: collapse to dim~$1$]\label{prop:step1}
$\widetilde W_1 = \mathrm{span}(u^+_{(2,n-1)})$, $1$-dimensional.
\end{proposition}

\begin{proof}
By Corollary~\ref{cor:Rn-image}, $W_n$ has the basis
$\{u^+_{(c,n-1)} : 2\le c\le n-2\}\cup\{v_{(n-1,n)}\}$.  For $c\ge 3$:
$u^+_{(c,n-1)}=\alpha v_{(c,n-1)}+\beta v_{(c,n)}$ has supports with
$2\notin S$ (since $S=(c,n-1)$ or $(c,n)$ with $c\ge 3$).  By
Remark~\ref{rem:E1-hook}, $u^+_{(c,n-1)}\in E_1^-$, so
$(T_1{+}1)u^+_{(c,n-1)}=0$.  Similarly $v_{(n-1,n)}\in E_1^-$ for
$n-1\ge 3$ (i.e.\ $n\ge 4$).

Only $u^+_{(2,n-1)}=\alpha v_{(2,n-1)}+\beta v_{(2,n)}$ survives, and
$\chi(2)=1$ for both supports, so $u^+_{(2,n-1)}\in E_1^+$ and
$(T_1{+}1)u^+_{(2,n-1)}=(1{+}q)u^+_{(2,n-1)}$.
\end{proof}

\begin{proposition}[Steps $2$ through $n-3$: shifting window]\label{prop:shift}
For $2\le j\le n-3$,
\[
   \widetilde W_j = \mathrm{span}(\widetilde u_j),\qquad
   \widetilde u_j \;=\; A_j v_{(j,n-1)}+B_j v_{(j+1,n-1)}+C_j v_{(j,n)}+D_j v_{(j+1,n)},
\]
with $A_j,B_j,C_j,D_j\in\mathbb{Q}(q)$ all nonzero (depending on $j,n,q$).
\end{proposition}

\begin{proof}
By induction on $j$.

\textbf{Base case $j=2$.}  Apply $(T_2{+}1)$ to $\widetilde W_1=
\mathrm{span}(u^+_{(2,n-1)})$.  By the analysis in
Proposition~\ref{prop:Wj} applied at $i=2$:
\begin{itemize}[leftmargin=2em,topsep=0.3em,itemsep=0.2em]
\item $T_2v_{(2,n-1)}$: $2$ at $(1,2)$, $3$ at $(2,1)$ (column-$1$
position).  $T_2$-block with $v_{(3,n-1)}$, axial distance $d=2$.
$(T_2{+}1)v_{(2,n-1)}=c_2 u^+$ where $u^+=\alpha_2 v_{(2,n-1)}+\beta_2 v_{(3,n-1)}$.

\item $T_2v_{(2,n)}$: similarly $T_2$-block with $v_{(3,n)}$ at $d=2$.
$(T_2{+}1)v_{(2,n)}=c_2(\alpha_2 v_{(2,n)}+\beta_2 v_{(3,n)})$.
\end{itemize}
So $(T_2{+}1)u^+_{(2,n-1)} = (1{+}q)c_2(\alpha\alpha_2 v_{(2,n-1)} +
\alpha\beta_2 v_{(3,n-1)} + \beta\alpha_2 v_{(2,n)} + \beta\beta_2 v_{(3,n)})$
where $(\alpha,\beta)=(\alpha_{n-1},\beta_{n-1})$ (the $T_{n-1}$-block
coefficients from $u^+_{(2,n-1)}$).  All four coefficients nonzero.

\textbf{Inductive step.}  Assume $\widetilde W_{j-1}$ is the span of
$\widetilde u_{j-1}$ supported on
$\{(j-1,n-1),(j,n-1),(j-1,n),(j,n)\}$ with all coefficients nonzero.

Apply $(T_j{+}1)$.  By Lemma~\ref{lem:Ti-hook}:
\begin{itemize}[leftmargin=2em,topsep=0.3em,itemsep=0.2em]
\item $v_{(j-1,n-1)}$ and $v_{(j-1,n)}$: $S=(j-1,n-1)$ resp.\ $(j-1,n)$
contain $j-1$ but not $j$ (since $j\ne n-1,n$ for $j\le n-3$).  In
column~$1$, $j$ is at cell $(j-1,1)$ and $j+1$ at $(j,1)$ (same column,
adjacent --- by counting as in Prop.~\ref{prop:Wj} step (i)).  So
$T_j$ acts as $-1$, and $(T_j{+}1)$ kills these.

\item $v_{(j,n-1)}$ and $v_{(j,n)}$: $S$ contains $j$ at $(1,2)$ but
not $j+1$ (since $j+1\ne n-1,n$ for $j\le n-3$).  $T_j$-block:
$\{v_{(j,b)},v_{(j+1,b)}\}$ at axial distance $d=j$ (for $b=n-1$ or $n$).
$(T_j{+}1)v_{(j,b)} = c_j(\alpha_j v_{(j,b)}+\beta_j v_{(j+1,b)})$.
\end{itemize}
Hence
\[
   (T_j{+}1)\widetilde u_{j-1} = c_j\bigl(B_{j-1}(\alpha_j v_{(j,n-1)}+\beta_j v_{(j+1,n-1)}) +
   D_{j-1}(\alpha_j v_{(j,n)}+\beta_j v_{(j+1,n)})\bigr),
\]
which is a single nonzero vector supported on $\{(j,n-1),(j+1,n-1),
(j,n),(j+1,n)\}$ with all coefficients nonzero (since $B_{j-1},D_{j-1},
\alpha_j,\beta_j\ne 0$).
\end{proof}

\begin{proposition}[Step $n-2$: terminal shape]\label{prop:terminal}
$\widetilde W_{n-2}$ is $1$-dimensional, supported on the $5$ subsets
\[
   \{(n{-}3,n{-}2),\;(n{-}3,n{-}1),\;(n{-}2,n{-}1),\;(n{-}2,n),\;(n{-}1,n)\},
\]
and contained in $E_1^-$.
\end{proposition}

\begin{proof}
Apply $(T_{n-2}{+}1)$ to $\widetilde W_{n-3}$, which is the span of
$\widetilde u_{n-3}$ supported on $\{(n{-}3,n{-}1),(n{-}2,n{-}1),
(n{-}3,n),(n{-}2,n)\}$.

Compute $(T_{n-2}{+}1)$ on each support:
\begin{itemize}[leftmargin=2em,topsep=0.3em,itemsep=0.2em]
\item $v_{(n-3,n-1)}$: $S=\{n{-}3,n{-}1\}$.  $n-2\notin S$ in column~$1$;
$n-1\in S$ at $(1,3)$.  Column~$1$ set sorted: $\{2,\ldots,n-4,n-2,n\}$
(omitting $n-3,n-1$).  Position of $n-2$: it is the $(n-4)$th smallest
(below it: $\{2,\ldots,n-4\}$, $n-5$ elements), at cell $(n-3,1)$,
content $-(n-4)$.  $T_{n-2}$-block: swap $n-2,n-1$ gives new SYT
$T_{(n-3,n-2)}$ ($n-1$ moves to column~$1$, $n-2$ moves to $(1,3)$).
Axial distance $d=2-(-(n-4))=n-2$.  Hence
$(T_{n-2}{+}1)v_{(n-3,n-1)}=c\cdot(\alpha v_{(n-3,n-1)}+\beta v_{(n-3,n-2)})$.
\item $v_{(n-2,n-1)}$: $S=\{n{-}2,n{-}1\}$, both at $(1,2),(1,3)$.
Same row.  $(T_{n-2}{+}1)v=(1{+}q)v_{(n-2,n-1)}$.
\item $v_{(n-3,n)}$: $S=\{n{-}3,n\}$.  $n-2,n-1\notin S$, both in
column~$1$.  Column~$1$ set: $\{2,\ldots,n-4,n-2,n-1\}$.  $n-2$ at
$(n-3,1)$, $n-1$ at $(n-2,1)$.  Same column, adjacent.
$(T_{n-2}{+}1)v=0$.
\item $v_{(n-2,n)}$: $S=\{n{-}2,n\}$.  $n-2$ at $(1,2)$.  $n-1\notin S$
in column~$1$: column-$1$ set $\{2,\ldots,n-3,n-1\}$, $n-1$ at last
position $(n-2,1)$, content $-(n-3)$.  $T_{n-2}$-block: swap
$n-2,n-1$ gives $T_{(n-1,n)}$.  Axial distance $d=1-(-(n-3))=n-2$.
$(T_{n-2}{+}1)v_{(n-2,n)}=c\cdot(\alpha v_{(n-2,n)}+\beta v_{(n-1,n)})$.
\end{itemize}

So if $\widetilde u_{n-3} = A v_{(n-3,n-1)} + B v_{(n-2,n-1)} +
C v_{(n-3,n)} + D v_{(n-2,n)}$ with $A,B,C,D\ne 0$, then
\[
   \widetilde u_{n-2} := (T_{n-2}{+}1)\widetilde u_{n-3} =
   c\alpha A\, v_{(n-3,n-1)} + c\beta A\, v_{(n-3,n-2)} +
   B(1{+}q)\, v_{(n-2,n-1)} + 0 + c\alpha D\, v_{(n-2,n)} + c\beta D\, v_{(n-1,n)}.
\]
Five nonzero coefficients; supports as claimed.

For $E_1^-$ containment: every support has minimum element $\ge n-3$,
which is $\ge 3$ for $n\ge 6$.  Hence $2\notin S$ for every support,
so by Remark~\ref{rem:E1-hook}, $\widetilde u_{n-2}\in E_1^-$.
\end{proof}

\begin{proof}[Proof of Theorem~\ref{thm:hook3}]
Combine Corollary~\ref{cor:Rn-image} (Phase~A) and
Propositions~\ref{prop:step1}, \ref{prop:shift}, \ref{prop:terminal}
(Phase~B).  $B_{n-1}V_\lambda=R'_{n-1}R'_n V_\lambda=\widetilde W_{n-2}$
is one-dimensional, supported on the $5$ subsets listed, and contained
in $E_1^-$.
\end{proof}

\section{Computational verification}

\begin{theorem}[Verification at $n=6,7,8$]\label{thm:verify}
Theorem~\ref{thm:hook3} and the intermediate steps
(Propositions~\ref{prop:Wj}, \ref{prop:step1}, \ref{prop:shift},
\ref{prop:terminal}) hold computationally at $n=6,7,8$.
\end{theorem}

\begin{proof}[Verification]
SymPy at $q=7\in\mathbb{Q}$ (generic for the seminormal basis).  At
each step, the rank, support pattern, and $E_1^-$ containment match
the predicted values.  At $n=8$, $\widetilde W_{n-2}$ is supported on
$\{(5,6),(5,7),(6,7),(6,8),(7,8)\}$, all with min element~$5\ge 3$,
hence in $E_1^-$.  Script:
\texttt{\textasciitilde/projects/scratch/2026-05-12-hook-j-formula/verify\_k3\_proof.py}.
\end{proof}

\section{The Catalan support phenomenon}

The structural proof above suggests a much broader phenomenon for
arbitrary hooks $\lambda=(k,1^{n-k})$.  Computational evidence at
$n\le 8$ for $k=2,3,4$ supports the following.

\begin{conjecture}[Catalan support for hooks]\label{conj:catalan}
Let $\lambda=(k,1^{n-k})\vdash n$ be a hook in the rank-zero regime
($\ell(\lambda)=n-k+1>\lceil n/2\rceil$, equivalently $n\ge 2k$).
Then the unique (up to scalar) spanning vector of
$B_{\ell+1}V_\lambda$ has support
\begin{enumerate}[leftmargin=2em,topsep=0.3em,itemsep=0.2em,label=(\arabic*)]
\item contained in subsets of the ``top range''
$R_k(n) := \{n-2k+3,\,n-2k+4,\,\ldots,\,n\}$ (size $2k-2$);
\item of cardinality exactly $C_k$, the $k$-th Catalan number.
\end{enumerate}
\end{conjecture}

\begin{theorem}[Computational verification of Conjecture~\ref{conj:catalan}]
\label{thm:catalan-verify}
Conjecture~\ref{conj:catalan} holds for $(k,n)\in\{(2,n) : n\ge 4\}\cup
\{(3,n) : n\ge 6\}\cup\{(4,8)\}$ at $n\le 8$.
\end{theorem}

\begin{proof}[Verification]
Compute $B_{\ell+1}V_\lambda$ in the seminormal basis at $q=7$,
identify the support of any nonzero column.  Results in
Table~\ref{tab:catalan}.

\begin{center}
\begin{table}[h]
\centering
\renewcommand{\arraystretch}{1.05}
\begin{tabular}{l l l c c c}
\toprule
$k$ & $n$ & $\lambda$ & top range $R_k(n)$ & $|$supp$|$ & $C_k$ \\
\midrule
2 & 4 & $(2,1,1)$            & $\{3,4\}$         & 2 & 2 \\
2 & 5 & $(2,1^3)$            & $\{4,5\}$         & 2 & 2 \\
2 & 6 & $(2,1^4)$            & $\{5,6\}$         & 2 & 2 \\
2 & 7 & $(2,1^5)$            & $\{6,7\}$         & 2 & 2 \\
3 & 6 & $(3,1,1,1)$          & $\{3,4,5,6\}$     & 5 & 5 \\
3 & 7 & $(3,1^4)$            & $\{4,5,6,7\}$     & 5 & 5 \\
3 & 8 & $(3,1^5)$            & $\{5,6,7,8\}$     & 5 & 5 \\
4 & 8 & $(4,1,1,1,1)$        & $\{3,4,5,6,7,8\}$ & 14 & 14 \\
\bottomrule
\end{tabular}
\caption{Verification of the Catalan support conjecture.}
\label{tab:catalan}
\end{table}
\end{center}

In every case, the support lies within the predicted top range and has
cardinality exactly $C_k$.  Script:
\texttt{\textasciitilde/projects/scratch/2026-05-12-hook-j-formula/extract\_support.py}.
\end{proof}

\subsection{Examples of the Catalan support}

\begin{example}[$k=3,n=6$]
The $\binom{4}{2}=6$ pairs in $R_3(6)=\{3,4,5,6\}$ are
$(3,4),(3,5),(3,6),(4,5),(4,6),(5,6)$.  The support of $B_5V_{(3,1,1,1)}$
is the $C_3=5$ pairs $\{(3,4),(3,5),(4,5),(4,6),(5,6)\}$ ---
all except the ``extreme'' pair $(3,6)$.
\end{example}

\begin{example}[$k=4,n=8$]
The $\binom{6}{3}=20$ triples in $R_4(8)=\{3,\ldots,8\}$ split as
$C_4=14$ supported and $\binom{2k-2}{k-3}=\binom{6}{1}=6$ excluded.
The $6$ excluded triples are
\[
   \{(3,4,7),(3,4,8),(3,5,8),(3,6,8),(3,7,8),(4,7,8)\}.
\]
A combinatorial characterization of the $C_k$-element family remains
to be identified (see Section~\ref{sec:open}).
\end{example}

\begin{remark}[Catalan identity]
$C_k = \binom{2k-2}{k-1}-\binom{2k-2}{k-3}$ is a standard form of the
Catalan number (e.g.\ via the ballot problem).  For $k=4$:
$\binom{6}{3}-\binom{6}{1}=20-6=14=C_4$.  For $k=5$:
$\binom{8}{4}-\binom{8}{2}=70-28=42=C_5$.
\end{remark}

\subsection{Implications for the $j$-formula}

\begin{corollary}[Conditional]\label{cor:j-from-catalan}
If Conjecture~\ref{conj:catalan} holds for $\lambda=(k,1^{n-k})$ in the
rank-zero regime, then the $j$-formula
(\texttt{2026-05-11-j-formula.tex}, Conj.~1) holds for $\lambda$, and
moreover the rank-corner formula gives the correct rank
($\kappa(\lambda)=1$).
\end{corollary}

\begin{proof}
The top range $R_k(n)=\{n-2k+3,\ldots,n\}$ has minimum element
$n-2k+3$.  For $\lambda$ in the rank-zero regime $n\ge 2k$, this
minimum is $\ge 3>2$, so every supported subset $S$ has $2\notin S$,
hence $v_S\in E_1^-$ by Remark~\ref{rem:E1-hook}.  The support has
$C_k\ne 0$ elements, hence the spanning vector is nonzero, giving
rank exactly $1$.
\end{proof}

\section{Open questions}\label{sec:open}

\begin{enumerate}[leftmargin=2em,topsep=0.3em,itemsep=0.2em]
\item \textbf{Combinatorial characterization of the Catalan support
  family.}  Identify a natural rule (Dyck paths, lattice paths,
  ballot sequences, etc.)\ describing which $C_k$ subsets of the top
  range $R_k(n)$ form the support.

\item \textbf{Structural proof of Conjecture~\ref{conj:catalan} for
  $k\ge 4$.}  The argument for $k=3$ in this paper relies on the image
  remaining $1$-dimensional throughout the second staircase factor
  $R'_{n-1}$.  For $k\ge 4$, the image of $B_{\ell+1}V_\lambda$ is still
  $1$-dimensional, but intermediate images $B_jV_\lambda$ for
  $\ell+2\le j\le n$ have higher dimension (specifically
  $\binom{n-1-(n-j+1)}{k-1-(n-j+1)}$).  A multi-dimensional shifting
  argument is needed.

\item \textbf{Beyond hooks.}  For non-hook $\lambda$ in the rank-zero
  regime (e.g.\ $\lambda=(2,2,1^{n-4})$), the Catalan support pattern
  presumably generalizes, with the support governed by a
  $\kappa(\lambda)$-tuple of Catalan-like structures (one per
  length-preserving corner).  Verifying at $n\le 8$ would be the
  natural next step.

\item \textbf{Closing the rank-zero theorem at $n\ge 8$ even.}  The
  May-11 j-formula (Cor.~1) reduces the rank-zero theorem to the
  $j$-formula at boundary $n=2k$.  Conjecture~\ref{conj:catalan} for
  $k=4,5,\ldots$ would close the boundary cases successively.
\end{enumerate}

\section{Summary}

\paragraph{Proven structurally.}
\begin{enumerate}[leftmargin=2em,topsep=0.3em,itemsep=0.2em]
\item The $j$-formula for $\lambda=(3,1^{n-3})$ at every $n\ge 6$
(Theorem~\ref{thm:hook3}).
\item Hence $\Pi^{S_n}|_{V_{(3,1^{n-3})}}=0$ unconditionally for every
$n\ge 6$ (Corollary~\ref{cor:hook3-rank-zero}).  Combined with the
May-11 result for $(2,1^{n-2})$, the rank-zero theorem now has a
structural proof for two infinite families of hooks.
\end{enumerate}

\paragraph{Computationally verified, structurally open.}
\begin{enumerate}[leftmargin=2em,topsep=0.3em,itemsep=0.2em]
\item The Catalan support pattern (Conjecture~\ref{conj:catalan}) at
$n\le 8$, $k\le 4$ (Theorem~\ref{thm:catalan-verify}).
\end{enumerate}

\paragraph{Why the $k=3$ proof works.}  Three structural facts conspire:
\begin{itemize}[leftmargin=2em,topsep=0.3em,itemsep=0.2em]
\item After $R'_n$, the image $W_n$ is the span of $+q$-eigenvectors of
$T_{n-1}$ at axial distance $n-1$, supported on subsets with second
element $\in\{n-1,n\}$ (Cor.~\ref{cor:Rn-image}).
\item The first factor $(T_1{+}1)$ of $R'_{n-1}$ collapses $W_n$ to the
$1$-dim subspace $\mathrm{span}(u^+_{(2,n-1)})$
(Prop.~\ref{prop:step1}).
\item Subsequent factors $(T_2{+}1),\ldots,(T_{n-2}{+}1)$ shift this
$1$-dim subspace through the SYT lattice in a controlled way: the
\emph{shifting window} of Prop.~\ref{prop:shift} takes the support
from $a\in\{2,3\}$ to $a\in\{n-3,n-2\}$, eventually landing in the
$5$-element configuration of Prop.~\ref{prop:terminal}.
\end{itemize}
For $k\ge 4$, the second step does not collapse to $1$-dim; instead the
image of $(T_1{+}1)\cdot W_n$ has dimension
$\binom{n-3}{k-3}>1$.  A higher-dimensional shifting argument is
required.

\paragraph{Honest assessment.}  This paper closes the rank-zero theorem
for one more infinite family ($k=3$ hooks) and provides strong
computational evidence for a Catalan structure governing the kernel of
$B_{\ell+1}$ on all hook irreducibles.  The structural proofs for
$k\ge 4$, the combinatorial characterization of the $C_k$-family, and
the boundary $n=2k$ cases for the rank-zero theorem all remain open.

\end{document}
